\section{Introduction}

Document-level relation extraction (DocRE) requires identifying semantic relations between entity pairs across multiple sentences, a task where encoder-based methods such as \citet{zhou2021atlop} long dominated.
Recent LLM-based approaches have closed this gap: \citet{autore2024} fine-tune LLMs with structured decomposition, \citet{xu2024lmrc} frame it as multi-choice reading comprehension, and \citet{chen2025docksrag} combine retrieval-augmented generation with structured knowledge.
These methods approach or surpass encoder baselines on DocRED \citep{yao2019docred} and Re-DocRED \citep{tan2022redocred} in extraction accuracy, yet every one outputs bare relation triplets $(e_h, r, e_t)$ without indicating \emph{which sentences} support each prediction.
In deployment contexts---clinical knowledge bases, legal due-diligence pipelines, financial compliance systems---users cannot verify a prediction without manually locating its textual basis across a multi-page document.

Encoder-based methods have partially addressed this need: \citet{xiao2022sais} introduce sentence-level evidence supervision as an auxiliary signal, and \citet{ma2023dreeam} extend this with dynamic evidence retrieval through memory attention.
Both produce evidence and improve extraction, but encoder methods now trail LLM-based approaches in accuracy.
The community faces a gap: the most accurate extractors provide no evidence, while evidence-capable methods are no longer competitive.

This gap reflects a structural difference.
In encoder methods, evidence retrieval is classification that regularizes the encoder at negligible cost; in LLM-based DocRE, evidence requires the model to \emph{generate} additional output tokens per triple, increasing output length and decoding complexity.
The na\"{i}ve expectation is that this generation cost always harms extraction accuracy.
Our experiments show this expectation is wrong.
The evidence tax is not a fixed penalty but a landscape shaped by model capacity, optimization paradigm, and training data quality.
Within the Qwen3 family, evidence generation \emph{benefits} extraction at 1.7B ($-$5.63pp Ign-F1, partly through halving output duplication) but is approximately neutral at $\geq$4B (+0.25$\pm$0.47pp at 4B).
Gradient-based optimization amplifies this competition---GRPO collapses extraction under constrained precision, DPO progressively trades extraction for evidence---while discrete sample selection (RSFT) avoids the trade-off entirely by filtering outputs rather than pushing gradients.

We formalize evidence-grounded DocRE as a structured generation task and compare three post-training paradigms---GRPO \citep{shao2024deepseekmath}, DPO \citep{rafailov2023direct}, and RSFT (building on STaR; \citealt{zelikman2022star})---across three fully evaluated scales (1.7B, 4B, 8B) with 14B validation, two model families, and two datasets. Our contributions:
\begin{enumerate}
    \item \textbf{Evidence generation is not uniformly harmful.} We observe a scale-dependent evidence tax with a within-family capacity threshold: beneficial at 1.7B ($-$5.63pp) but approximately neutral at $\geq$4B (+0.25$\pm$0.47pp at 4B), challenging the assumption that evidence always costs extraction accuracy.

    \item \textbf{Gradient pressure causes evidence-extraction decoupling.} We characterize a decoupling spectrum (GRPO catastrophic $\rightarrow$ DPO gradual $\rightarrow$ RSFT none) and introduce the Evidence-Decoupling Ratio (EDCR) as a diagnostic, validated at per-relation level ($\rho{=}{-}0.79$, $n{=}76$, $p{<}0.001$).

    \item \textbf{RSFT consistently reduces training variance.} NLI-based sample selection improves mean Ign-F1 across tested scales (+1.2 to +6.6pp); our 1.7B RSFT model (51.19 Ign-F1, std 0.09pp) approaches 8B SFT performance (50.19, std 3.23pp) at 4.7$\times$ fewer parameters, enabling deployment-efficient evidence-grounded extraction.
\end{enumerate}
